\documentclass[10pt,a4paper]{amsart}
\usepackage[margin=19mm]{geometry}
\usepackage{amsmath,amssymb,mathtools}
\usepackage[scr=rsfs,cal=euler]{mathalfa}
\usepackage{xfrac}
\usepackage[unicode,hidelinks,bookmarks=false]{hyperref}
\newcommand{\ie}{i.e.\ }
\newcommand{\cf}{cf.\ }
\newcommand{\resp}{respectively\ }
\newcommand{\Subsection}{\subsection}
\providecommand{\nobreakdash}{\nobreak\hbox{-}}
\setlength{\emergencystretch}{2em}
\multlinegap0pt
\usepackage{bm}
\usepackage{bbm}
\usepackage{dsfont}
\usepackage{xspace}
\usepackage{cancel}
\usepackage{mathtools}
\usepackage{amsthm}
\makeatletter
\@ifundefined{thm}{\newtheorem{thm}{Thm}}{}
\@ifundefined{lemma}{\newtheorem{lemma}[thm]{Lemma}}{}
\@ifundefined{theorem}{\newtheorem{theorem}[thm]{Theorem}}{}
\makeatother
\makeatletter
\newcommand{\Gal}{{\rm Gal}}
\expandafter\ifx\csname acknowledgements\endcsname\relax
\newenvironment{acknowledgements}
{\par\noindent\textbf{Acknowledgements.}\ }
\fi
\makeatother
\title[On the Vanishing of the Brauer-Manin Obstruction for Normic ]{On the Vanishing of the Brauer-Manin Obstruction for Normic Bundles}
\author{Mridul Biswas, Divyasree C-Ramachandran, Biswanath Samanta}
\date{Source: arXiv:2607.25287v1 (CC BY 4.0)}
\begin{document}
\maketitle

\noindent\emph{Source and licence.} On the Vanishing of the Brauer-Manin Obstruction for Normic Bundles. Version arXiv:2607.25287v1 (CC BY 4.0), distributed under the Creative Commons Attribution 4.0 International licence, as recorded in the arXiv licence metadata for this exact version and shown on the abstract page https://arxiv.org/abs/2607.25287v1. Full rights notice: licenses/arxiv-2607.25287-CC-BY-4.0.txt.\par\smallskip

\noindent\textit{Editorial scope.} Counted unit: the Lemma whose source label is \texttt{field-theory lemma} in the article (source0.tex lines 471--477, source label \texttt{field-theory lemma}), delivered together with the article's own complete original proof of it (source0.tex lines 479--506, one \texttt{proof} environment of 28 lines). No mathematics is written, repaired or re-derived: statement and proof are the article's own text line for line. Required context retained uncounted: Frattini's argument (lines 449--451); Schur--Zassenhaus (lines 456--458). Every definition, display and label of those lines is retained unchanged. Omitted from this package and not redistributed: 1--448, 452--455, 459--470, 478--478, 507--925. Nothing inside a retained excerpt is omitted or altered: every retained body is delivered exactly as the article's lines are written, so no omission receipt of any kind accompanies this package. Citations: the retained excerpts carry no citation command at all, so no bibliography entry of the article is transcribed and no reference is redistributed. Reference adaptation: none was needed and none was made; every reference inside the delivered unit resolves inside it, so no unresolved reference or citation is produced. Printed slips kept literal: no printed word, symbol, hypothesis or number of the counted statement or of its proof was altered. Layout and class: the article's own class is \texttt{amsart} with packages including geometry, amssymb, amsmath, amsthm, amstext, amscd, latexsym, graphics, graphicx, bbm, caption, xcolor, hyperref, mathtools; the wrapper is the lane's pinned \texttt{amsart} editorial wrapper at 10pt on A4, which restates the theorem-like environments the retained text needs and adds the article's own macro definitions that the retained text needs (\texttt{\textbackslash Gal}), with extra wrapper packages (bm, bbm, dsfont, xspace, cancel, mathtools). The delivered numbering is the unit's own. Assets: no retained span contains a figure environment, a graphics-inclusion command, a PDF-page inclusion, a table or a TikZ picture; no image file is copied into this package, the package declares no asset dependency and no image file is redistributed with it. Licence: version arXiv:2607.25287v1 of this article is distributed under the Creative Commons Attribution 4.0 International licence, as recorded in the arXiv licence metadata for this exact version and shown on the abstract page https://arxiv.org/abs/2607.25287v1. A full rights notice is delivered at licenses/arxiv-2607.25287-CC-BY-4.0.txt. Characters: every retained excerpt is pure ASCII, so the delivered engine, XeLaTeX with its default Latin Modern text font, reports no missing character.
\begin{lemma}[Frattini]\label{Frattini's argument}
Let $G$ be a finite group and $H$ be a normal subgroup of $G$. If $P$ is a Sylow $p$-subgroup of $H$, then $G=N_G(P)\,H$.
\end{lemma}

\begin{theorem}[Schur--Zassenhaus]\label{Schur--Zassenhaus}
Let $|G|=ab$ with $(a,b)=1$. If $G$ has a normal subgroup $H$ of order $a$, then $G$ has a subgroup $S$ of order $b$, and $G=HS\cong H\rtimes S$.
\end{theorem}

\begin{lemma}\label{field-theory lemma}
Let \(L/k\) be a finite Galois extension such that \(p\mid [L:k]\). Let
\(E_1,\dots,E_r\) be Galois subfields of \(L/k\) satisfying
\(p\nmid [E_i:k]\) for each \(i\). Then there exists an intermediate field
\(R\) of \(L/k\) such that \(p\mid [R:k]\) and \(R\cap E_i=k\) for all
\(1\le i\le r\).
\end{lemma}

\begin{proof}
Let \(G=\Gal(L/k)\), and let \(C\) be the compositum of the fields
\(E_1,\dots,E_r\). Since each extension \(E_i/k\) has degree coprime to \(p\),
we have \(p\nmid [C:k]\). Let \(H\trianglelefteq G\) be the subgroup corresponding to \(C\).

Let \(P\) be a Sylow \(p\)-subgroup of \(G\). Since \(p\nmid [G:H]\), we have
\(P\le H\). By Theorem~\ref{Schur--Zassenhaus}, there exists a subgroup
\(S\le N_G(P)\) such that \(N_G(P)=P\rtimes S\). In particular, \(P\cap S=\{1\}\), and hence
\(p\nmid |S|\). Since \(p\mid |G|\), it follows that \(p\mid [G:S]\). As \(P\le H\trianglelefteq G\), Lemma~\ref{Frattini's argument} yields
\(G=N_G(P)H\). Using \(N_G(P)=PS\) and \(P\le H\), we obtain \(G=(PS)H=SH\).

Set \(R=L^S\). Then \(p\mid [R:k]\). For each \(i\), let \(H_i\le G\) be the
subgroup corresponding to \(E_i\). Since \(E_i\subseteq C\), we have
\(H_i\supseteq H\). Therefore,
\(SH_i\supseteq SH=G\), and hence \(SH_i=G\). By the Galois correspondence,
$$
R\cap E_i
=
L^S\cap L^{H_i}
=
L^{SH_i}
=
L^G
=
k.
$$
Thus \(R\cap E_i=k\) for all \(1 \le i \le r\).
\end{proof}

\end{document}
